% =====================================================================
% Module 03 — Spectral Bundle Structure
% Academic-Galactic v1.2
% =====================================================================

\section{Spectral Bundle Structure}
\label{sec:SpectralBundleStructure}

\subsection{Overview}
Spectral bundle structure defines how spectral operators $\Phi_5$ act across
filtration layers and fiber bundles of the Soft 80-dimensional manifold
$\mathcal{M}^{80}$. Each filtration layer $\mathcal{F}_k$ carries a spectral
bundle
\[
\sigma_k : S_k \to \mathcal{F}_k,
\]
where $S_k$ encodes spectral curvature, spectral invariants, and operator
spectra.

Spectral geometry foundations follow operator-algebraic state space geometry
\citeClassA{AlfsenSchultz2003} and modern spectral curvature research
\citeClassA{SpectralCurvature2026}.

\subsection{Spectral Operators}
\begin{definition}[Spectral Operator]
A spectral operator at altitude $k$ is a map
\[
\Phi_5 : \mathcal{F}_k \to S_k
\]
assigning to each point its spectral signature.
\end{definition}

Spectral operators are compatible with fiber bundle geometry (Module 02) and
filtration structure (Module 01).

\subsection{Spectral Curvature}
Spectral curvature is defined as the variation of spectral signatures along
curves $\gamma : [0,1] \to \mathcal{F}_k$:
\[
\kappa_{\mathrm{spec}}(\gamma)
=
\left\| \frac{d}{dt} \Phi_5(\gamma(t)) \right\|.
\]

Spectral curvature satisfies nonnegativity constraints
\citeClassA{SpectralCurvature2026} and interacts with curvature envelopes
$\Phi_7$ (Module 02).

\subsection{Spectral Continuity}
Spectral continuity ensures that spectral signatures vary smoothly across
filtration layers:
\[
\Phi_5(i_k(p)) = \Phi_5(p),
\]
where $i_k : \mathcal{F}_k \hookrightarrow \mathcal{F}_{k+1}$ is the filtration
embedding.

This follows from operator-algebraic continuity
\citeClassA{AlfsenSchultz2003}.

\subsection{Spectral Invariants}
Spectral invariants are quantities preserved under operator action:
\[
I_{\mathrm{spec}}(p) = I_{\mathrm{spec}}(\Phi_8(p)),
\]
where $\Phi_8$ is the holonomy operator (Module 02).

Examples include:
\begin{itemize}
    \item spectral radius,
    \item spectral gap,
    \item spectral curvature integral.
\end{itemize}

\subsection{Spectral–Dual Alignment}
Dual compatibility (Module 06) requires spectral–dual alignment:
\[
\mathcal{C}(\Phi_5(p)) = \Phi_5(\mathcal{C}(p)),
\]
where $\mathcal{C}$ is the dual operator.

This ensures:
\begin{itemize}
    \item cotangent spectral consistency,
    \item dual-operator stability,
    \item reprojection compatibility (Module 07).
\end{itemize}

\subsection{Spectral Bundle Local Triviality}
Each spectral bundle $S_k$ is locally trivial:
\[
\sigma_k^{-1}(U) \cong U \times \Lambda_k,
\]
where $\Lambda_k$ is the spectral fiber at altitude $k$.

Local triviality ensures compatibility with:
\begin{itemize}
    \item fiber bundle geometry (Module 02),
    \item collapse-surface stratification (Module 04),
    \item composite operator algebra (Module 05).
\end{itemize}

\subsection{Summary}
The spectral bundle structure defines spectral operators, spectral curvature,
spectral continuity, spectral invariants, and spectral–dual alignment. These
structures interact with fiber bundle geometry, collapse-surface stratification,
and the composite operator algebra, forming the spectral backbone of the A7
operator altitude.

% =====================================================================
% End Module 03
% =====================================================================
